\DocumentMetadata{
  pdfversion=2.0,
  pdfstandard=ua-2,
  lang=en-US,
  tagging=on,
  tagging-setup={math/setup=mathml-SE,tabsorder=structure}
}
\documentclass[11pt,letterpaper]{article}
\usepackage[margin=0.68in,includefoot]{geometry}
\usepackage{newcomputermodern}
\usepackage{amsmath,amscd,mathtools}
\usepackage{tikz,adjustbox}
\providecommand{\square}{\mdlgwhtsquare}
\usepackage[hidelinks]{hyperref}
\hypersetup{
  pdftitle={Probability PhD Exam, May 1989},
  pdfauthor={University of Florida Department of Mathematics},
  pdfsubject={Graduate mathematics examination},
  pdfdisplaydoctitle=true,
  bookmarksopen=true
}
\setcounter{secnumdepth}{0}
\setlength{\parindent}{0pt}
\setlength{\parskip}{0.28em}
\setlength{\emergencystretch}{3em}
\setlength{\leftmargini}{2.15em}
\setlength{\leftmarginii}{2.65em}
\setlength{\leftmarginiii}{3em}
\renewcommand{\labelenumi}{\textbf{\theenumi.}}
\pagestyle{empty}
\raggedbottom
\allowdisplaybreaks
\newenvironment{examproblems}[1]{%
  \begin{enumerate}%
  \setcounter{enumi}{#1}%
  \setlength{\topsep}{0.35em}%
  \setlength{\partopsep}{0pt}%
  \setlength{\itemsep}{0.48em}%
  \setlength{\parsep}{0.18em}%
}{\end{enumerate}}
\newenvironment{examsubparts}{%
  \begin{enumerate}%
  \setlength{\topsep}{0.18em}%
  \setlength{\partopsep}{0pt}%
  \setlength{\itemsep}{0.16em}%
  \setlength{\parsep}{0.1em}%
}{\end{enumerate}}
\newenvironment{examinstructions}{%
  \par\begingroup\itshape
}{%
  \par\endgroup\vspace{0.18em}
}
\makeatletter
\renewcommand\section{\@startsection{section}{1}{0pt}%
  {-1.25ex plus -.25ex minus -.1ex}%
  {0.55ex plus .1ex}%
  {\normalfont\large\bfseries}}
\renewcommand{\@maketitle}{%
  \newpage\null\vskip -1em%
  {\footnotesize\bfseries
    \href{https://www.ufl.edu/}{University of Florida}, \href{https://math.ufl.edu/}{Department of Mathematics}\par}%
  \vskip -0.5em%
  \tagstructbegin{tag=Title}%
    {\LARGE\bfseries\href{https://gma.math.ufl.edu/exams/probability-phd/}{Probability PhD Exam}, May 1989\par}%
  \tagstructend%
  \vskip 0.25em%
  \tagmcbegin{artifact}%
    \hrule height 1pt%
  \tagmcend%
  \vskip 1em%
}
\makeatother
\title{Probability PhD Exam, May 1989}
\author{}
\date{}
\begin{document}
\maketitle
\thispagestyle{empty}
\begin{examproblems}{0}
\item
Let \(X_n\) be independent identically distributed random variables and \(S_n=\sum_{i=1}^n X_i\). Suppose 
\[
\mathrm E\left[\sup_n\left|\frac{S_n}{n}\right|\right]<\infty.
\]
 Show that 
\[
\mathrm E\bigl[|X_1|\log^+|X_1|\bigr]<\infty,
\]
 where 
\[
\log^+x=\begin{cases}
\log x,&x\geq 1,\\
0,&\text{otherwise}.
\end{cases}
\]
 Hint: \(C_n=\prod_{j=1}^n \mathrm P[|X_j|\leq j]\to c>0\). If \(T=\inf\{n:|X_n|>n\}\), then 
\[
\sum_{n=1}^{\infty}\frac1n\int_{\{T=n\}}|X_n|\,dp
=\sum_{n=1}^{\infty}\frac{C_{n-1}}n\int_{\{|X_1|>n\}}|X_1|\,dp
\]
 and 
\[
\sum_{n=1}^{\infty}\frac1n\int_{\{T=n\}}|S_{n-1}|\,dp<\infty.
\]
\item
Let \(X_n\) be independent Poisson variables. Suppose \(\sum_{n=1}^{\infty}X_n<\infty\) a.s. Show that \(\sum_{n=1}^{\infty}\mathrm E[X_n]<\infty\).

\par
Hint. Borel–Cantelli.
\item
This problem has two parts.
\begin{examsubparts}
\item[\textnormal{a)}]
Define the notion of conditional expectation. In particular explain the notion \(\mathrm E[X\mid Y]\) where~\(X\) and~\(Y\) are random variables.
\item[\textnormal{b)}]
Suppose \(X,Y\) are random variables such that \(\mathrm E(X^2),\mathrm E(Y^2)<\infty\) and 
\[
\mathrm E[X\mid Y]=Y,
\qquad
\mathrm E[Y\mid X]=X.
\]
 Show that \(X=Y\) a.s.
\end{examsubparts}
\item
This problem has three parts.
\begin{examsubparts}
\item[\textnormal{a)}]
A sequence of characteristic functions on \(\mathbb R^1\) converging pointwise to a characteristic function does so uniformly on compact sets. Explain why - no need to give complete proofs.
\item[\textnormal{b)}]
Let \(b_n\in\mathbb R\) and~\(f\) a characteristic function not identically equal to~\(1\). Suppose \(f(b_nt)\) converges to a characteristic function. Show that \(\lim b_n=b\in\mathbb R\).
\item[\textnormal{c)}]
Let \(\phi(t)=\frac18[1+7e^{it}]\). Show that~\(\phi\) is a characteristic function but that \(|\phi|\) is not.

\par
Hint. \(|\phi|^2\) is the characteristic function of a measure concentrated on \(\{-1,0,1\}\). If \(|\phi|\) is the characteristic function of a measure~\(m\), then \(|\phi|^2\) is the characteristic function of \(m*m\).
\end{examsubparts}
\item
Let \(X_t\) be a standard Brownian motion.
\begin{examsubparts}
\item[\textnormal{(i)}]
Show that \((X_t)\), \((X_t^2-t)\), and \(e^{uX_t-\frac12u^2t}\) are all martingales.
\item[\textnormal{(ii)}]
Let \(T_a(w)=\inf\{t:X_t(w)=a\}\) for \(a>0\), the inf being defined to be \(+\infty\) if the set is empty. Using (i), compute \(\mathrm E[e^{-\lambda T_a}]\) for \(\lambda>0\), and show that \(T_a<\infty\) a.s. and \(\mathrm E[T_a]=\infty\).
\end{examsubparts}
\item
State and sketch the proof of the martingale convergence theorem.
\item
Using the martingale convergence theorem prove the following:

\par
Let \(F_n\) be an increasing sequence of~\(\sigma\)-fields, \(X_n\) a sequence of random variables dominated by an integrable random variable~\(Y\): \(|X_n|\leq Y\). Suppose \(X_n\to X\) a.s. Then 
\[
\mathrm E[X_n\mid F_n]\to \mathrm E[X\mid F_\infty]\quad\text{a.s.}
\]
 Hint. Put the 
\[
U_n=\sup_{m\geq n}|X_m-X|.
\]
 Then for \(m\geq n\), 
\[
\mathrm E[|X_m-X|\mid F_m]\leq \mathrm E[U_n\mid F_m].
\]
 
\[
\mathrm E\left\{\limsup_m\mathrm E[|X_m-X|\mid F_m]\right\}
\leq \mathrm E\{\mathrm E[U_n\mid F_\infty]\}=\mathrm E[U_n].
\]
\end{examproblems}
\end{document}
